
\section{A topology library built from finite certificates}
\label{sec:library}

This section lists the principal topological notions and gives an \ipra-compatible formulation.  The goal is not to force all topology into one definition, but to expose exactly which finite observations and positive certificates each theorem uses.

\subsection{Open information and continuity}

At scale $k$, a basic open is a vertex of $K_k$.  A finite union is a bitset.  Membership of a presented point is positive when its observation simplex meets the bitset in a certified way.  Arbitrary subsets of points are not internal objects.

A continuous map is represented by a modulus and carriers as in \cref{sec:towers}.  This formulation has three advantages:
\begin{enumerate}
  \item it is pointfree at the finite level;
  \item it permits multiple valid target cells at boundaries;
  \item composition is finite relational composition.
\end{enumerate}

For metric presentations, the familiar evaluator/modulus formulation induces the carrier formulation.  Conversely, if the tower is spatially complete and has a shrinking mesh, a coherent carrier map determines the usual continuous action on regular point names.  That converse is a semantic theorem about the chosen presentation, not a generic fact of first-order arithmetic.

\subsection{Compactness}

Abstract compactness says every open cover has a finite subcover.  Quantifying over arbitrary open covers is higher-type and often nonconstructive.  The usable data are instead:

\begin{definition}[Finite compactness contract]
A compactness contract supplies a primitive-recursive net $N(k)$, a finite cover $\mathcal U_k$, and a verifier proving that every represented point is observed in at least one member, with mesh or Lebesgue data sufficient for the intended theorem.
\end{definition}

If an application receives a finite candidate cover by basic opens, coverage reduces to a bounded statement about the net plus a supplied uniform-continuity or inclusion margin.  Bounded search extracts a subcover when the list is redundant.  The contract does not turn an arbitrary infinite enumeration of opens into a finite subcover without a bar, modulus, or stopping bound.

\subsection{Separation and apartness}

A Hausdorff-style effective presentation supplies a separation procedure:
\[
  \operatorname{Apart}(x,y,s)=1
  \Longrightarrow
  \text{disjoint basic neighborhoods of scale }s.
\]
For regular fast metric names, a witnessed lower distance bound provides this procedure.  Negated name equality alone does not.

At the tower level, two observation clouds are apart at scale $k$ when no simplex contains a vertex from each cloud, or when an explicit disjointness certificate is supplied.  Separation propagates under refinement by a bounded proof.

\subsection{Connectedness and paths}

There are several constructive strengths.

\begin{definition}[Finite chain connection]
Two vertices of $K_k$ are chain connected when they lie in the same component of the $1$-skeleton.  A certificate is a finite vertex path whose consecutive pairs are edges.
\end{definition}

A presented space is \emph{uniformly chain connected} when a fixed primitive-recursive procedure produces such paths between observation clouds at every requested scale.  This positive statement is stronger and more computational than merely denying a separation into two opens.

A represented path is a tower map from an interval tower.  Polygonal paths with rational vertices are especially convenient: evaluation, subdivision, and clearance tests are primitive recursive.  Homotopies are tower maps from a product with an interval tower.

\subsection{Products}

Given covers $\mathcal U_k$ and $\mathcal V_k$, use the product cover $U_i\times V_j$.  A finite family of product cells has nonempty intersection exactly when both coordinate families do, provided the factor intersection certificates are available.  Thus the product nerve is a finite combinatorial construction on pairs.  Refinement and map carriers are componentwise.

The construction proves that products of certified towers are certified towers.  For metric spaces, the max metric gives a simple mesh bound.  All dimension bounds, bitset encodings, and demand transformations are explicit.

\subsection{Subspaces and closed conditions}

A subspace cannot in general be formed merely from a negative predicate.  The tower needs a positive answer to ``which cells meet the subspace?''  Suitable inputs include:
\begin{itemize}
  \item a finite decidable subset of a finite space;
  \item an zero set with isolating boxes and sign/degree certificates;
  \item a compact set with a distance-name and positive intersection witnesses;
  \item a supplied finite subcomplex.
\end{itemize}
Cells with certified intersection become vertices; inherited intersection certificates define the subcomplex.  Ambiguous cells remain as clouds rather than being discarded.

\subsection{Quotients}

For finite spaces and finite complexes, quotients by a decidable equivalence relation are bounded constructions.  For represented spaces, a quotient presentation requires more:
\begin{enumerate}
  \item a fixed evaluator for the equivalence relation or orbit action;
  \item finite saturation of observation clouds;
  \item compatibility of refinement;
  \item separation data if a Hausdorff quotient is claimed.
\end{enumerate}
Without these fields, a quotient may destroy effective separation or make equality undecidable.  The framework therefore treats quotient quality as a certificate, not as an automatic theorem.

\subsection{Homology and cohomology}

At each scale, mod-$2$ homology is computed from finite boundary matrices.  Gaussian elimination is a bounded loop.  Its invariant records:
\begin{itemize}
  \item the processed pivot columns;
  \item preservation of row space under elementary operations;
  \item pivot distinctness;
  \item rank accounting.
\end{itemize}
The Betti formula is
\[
  \beta_q(K)=n_q-\operatorname{rank}\partial_q-
  \operatorname{rank}\partial_{q+1}.
\]

Cohomology uses transposed boundary matrices.  Cup products are finite formulas on simplex tables; associativity and the Leibniz rule can be proved by bounded index manipulations.  Integral torsion is added by checked Smith certificates as described after \cref{thm:prism}.

The tower does not force Betti numbers to stabilize.  Stabilization, pro-isomorphism, or identification with singular/Cech homology is a semantic theorem requiring hypotheses on the cover system.  McCord's finite-space results and classical nerve theorems indicate the relevant bridges \cite{McCord1966}, but each effective use must supply finite good-cover or acyclicity certificates.

\subsection{Fundamental groups}

For a finite complex, an edge loop is a finite list of vertices.  Elementary edge homotopies insert or delete backtracks and triangle moves.  A bounded sequence of such moves is a positive null-homotopy certificate.  Van Kampen presentations can be produced from a finite spanning tree and $2$-simplex relations.

What is not available uniformly is a decision procedure for the word problem in arbitrary finitely presented groups.  The \ipra\ interface is therefore asymmetric:
\begin{itemize}
  \item a supplied bounded homotopy is verifiable;
  \item bounded search can find one below a stated bound;
  \item absence of a bounded certificate is not a proof of non-null-homotopy;
  \item special classes with a proven normal-form algorithm may add a decision contract.
\end{itemize}

\subsection{Manifolds and atlases}

A compact effective $d$-manifold can be presented by a finite atlas at each scale:
\[
  (U_i,\phi_i,\psi_i,r_i),
\]
where $\phi_i$ and $\psi_i$ are fixed rational evaluators for a chart and its inverse on guarded domains, $r_i$ is a positive interior margin, and overlap transitions have explicit moduli.  Finite cover certificates prove that the shrunken chart domains cover the space.  Carriers handle points lying in several charts.

Local topological proofs reduce to rational boxes.  Differential topology adds rational derivative names, determinant gaps for inverse-function steps, and transversality margins.  Smooth partitions of unity require explicit bump-function approximants and denominator lower bounds.  None of these should be hidden inside the word ``smooth.''

For noncompact manifolds, add a primitive-recursive exhaustion by compact charted pieces.  A theorem at precision $k$ advertises the exhaustion stage it needs.  This is the same finite-demand discipline used throughout.

\subsection{Fixed-point theorems}

A Sperner labeling on a finite triangulation is a bounded object.  The parity proof yields an approximate fixed-point simplex.  Given a uniformly continuous map on a compact presented simplex, a mesh choice turns this into
\[
  \|f(x)-x\|<2^{-k}
\]
for an explicitly produced rational point $x$.

An exact fixed point is a different claim.  To construct a coherent name, one needs nested compatible approximate witnesses.  This can be obtained from an isolation/uniqueness modulus, from a relational viability dynamic program, or from an explicit compact choice principle available for the particular presentation.  Unrelated approximate fixed points at each scale do not automatically form one point.

\subsection{Dimension and local degree}

Dimension obstructions are particularly suitable for finite topology.  The repository's staggered-grid theorem is already a quantitative finite alternative.  Invariance of domain can be approached similarly: triangulate a rational box, represent a map by a sufficiently fine carrier, compute a local mod-$2$ degree on a boundary complex, and use an injectivity-separation modulus to prevent cancellation.  The finite degree calculation is within \ipra; producing the required uniform separation data is the substantive application hypothesis.

One must not silently infer unrestricted classical invariance theorems from the finite core.  Reverse-mathematical work shows that some unrestricted Brouwer invariance statements have weak-Koenig-lemma strength over standard base theories \cite{Kihara2020}.  The present goal is the maximal explicit, effectively presented version: fixed evaluators, positive moduli, and finite certificates.
